%% cv_awesomecv.tex — Dipayan Sarkar (Awesome-CV style)
\documentclass[11pt, a4paper]{awesome-cv}

\geometry{left=1.4cm, top=1.5cm, right=1.4cm, bottom=1.5cm, footskip=.5cm}

\colorlet{awesome}{awesome-skyblue}
\setbool{acvSectionColorHighlight}{true}

% Personal info
\name{Dipayan}{Sarkar}
\position{PhD Researcher{\enskip\cdotp\enskip}Bioinformatics \& Computational Biology}
\address{Siliguri, West Bengal, India}
\mobile{+91 9735490241}
\email{dipayansarkar26@gmail.com}
\homepage{dipayansarkar.com}
\github{Dipayan26}
\linkedin{dipayansarkar}
\googlescholar{scholar}{Google Scholar}
\photo[circle,noedge,left]{dpn_photo}

\begin{document}

\makecvheader[C]

\makecvfooter{\today}{Dipayan Sarkar~~~·~~~CV}{\thepage}

% ---------- Research Summary ----------
\cvsection{Research Summary}
\begin{cvparagraph}
Currently pursuing a PhD in bioinformatics and computational biology, focused on deep learning methods for biological sequence-based prediction and generative model development. Experience includes neural network training, transfer learning with pretrained protein language models, and analysis of large-scale biological datasets. Current work involves autoregressive generative modeling of biological sequences, with ongoing exploration of diffusion-based approaches for biological data.
\end{cvparagraph}

% ---------- Work Experience ----------
\cvsection{Work Experience}
\begin{cventries}
  \cventry
    {PhD Research Scholar}
    {University of North Bengal}
    {Siliguri, India}
    {2022 -- Present}
    {
      \begin{cvitems}
        \item Designed and implemented deep learning architectures for sequence-based protein--protein interaction prediction.
        \item Developed attention-based hybrid models integrating transfer learning from pretrained protein language models, trained and evaluated on large-scale biological sequence datasets.
        \item Developed and deployed web-based platforms for model inference and protein interaction network visualization.
      \end{cvitems}
    }
\end{cventries}

% ---------- Education ----------
\cvsection{Education}
\begin{cventries}
  \cventry
    {Ph.D. in Bioinformatics}
    {University of North Bengal}
    {India}
    {Expected March 2027}
    {Advisor: Dr. Chiranjib Sarkar \quad|\quad Focus: Deep learning-based frameworks for protein--protein interaction network prediction.}

  \cventry
    {M.Sc. in Botany}
    {University of North Bengal}
    {India}
    {2021}
    {Specialization: Biochemistry \quad|\quad 77.25\% (First Class)}

  \cventry
    {B.Sc. (Hons.) in Botany}
    {Ananda Chandra College, University of North Bengal}
    {India}
    {2019}
    {63.50\% (First Class)}

  \cventry
    {Higher Secondary (Class XII)}
    {Haldibari High School (H.S.), WBCHSE}
    {West Bengal}
    {2016}
    {80.00\% (First Class) \quad|\quad Subjects: Physics, Chemistry, Mathematics, Biology, English, Bengali}

  \cventry
    {Secondary (Class X)}
    {Haldibari High School (H.S.), WBBSE}
    {West Bengal}
    {2014}
    {91.58\% (First Class)}
\end{cventries}

% ---------- Publications ----------
\cvsection{Publications}
\begin{cvparagraph}
\begin{enumerate}
  \item Sarkar, D., \& Sarkar, C. (2025). AttnSeq-PPI: Enhancing protein-protein interaction network prediction using transfer learning-driven hybrid attention. \textit{Biochimica et Biophysica Acta (BBA)---Proteins and Proteomics}, 141102.
  \item Sarkar, C.; Sarkar, D.; Parsad, R.; Mishra, D. (2022). Package \texttt{EGRNi}. CRAN. \url{https://cran.r-project.org/package=EGRNi}
\end{enumerate}
\end{cvparagraph}

% ---------- Skills ----------
\cvsection{Skills}
\begin{cvskills}
  \cvskill{Programming}{Python (PyTorch, HuggingFace), R (CRAN package dev), C (basic)}
  \cvskill{Deep Learning}{Transformers, attention mechanisms, sequence modeling, transfer learning}
  \cvskill{Generative AI}{Autoregressive sequence generation, diffusion-based approaches, probabilistic modeling}
  \cvskill{Scalable Training}{GPU training, large-scale datasets, parallel workflows on HPC}
  \cvskill{Bioinformatics}{Sequence analysis, sequence-only PPI prediction, generative modeling}
  \cvskill{Data \& Viz}{NumPy, Pandas, Matplotlib, Seaborn}
  \cvskill{Deployment}{Docker, Flask, FastAPI, Linux server deployment}
\end{cvskills}

% ---------- Software & Tools ----------
\cvsection{Software \& Tools}
\begin{cventries}
  \cventry
    {Deep learning-based PPI prediction platform}
    {AttnSeq-PPI}
    {}{}
    {\url{https://compbiosysnbu.in/attnseqppi/}}

  \cventry
    {Arabidopsis-specific PPI prediction and network analysis tool}
    {ARACoFusion-PPI}
    {}{}
    {\url{https://aracofusion.compbiosysnbu.in/}}
\end{cventries}

% ---------- Awards ----------
\cvsection{Awards \& Certificates}
\begin{cvhonors}
  \cvhonor
    {CSIR-NET (JRF)}
    {Life Sciences, All India Rank: 216}
    {India}
    {June 2021}

  \cvhonor
    {GATE}
    {Life Sciences (XL)}
    {India}
    {2022}
\end{cvhonors}

\end{document}
